\section{Introduction}
\label{sec:introduction}

Robotic construction is increasingly planned from digital building
information, yet the model delivered by a designer and the world required by a
robot remain fundamentally different artefacts. An Industry Foundation Classes
(IFC) model describes intended components, spatial containment, materials, and
design relationships. A robot simulator instead requires collision geometry,
mass properties, physical constraints, task targets, and an explicit account of
which objects and geometric details are active. The difference is not resolved
by changing file formats. It is a change from a representation of what should
be built to a representation of what can be acted upon, under a particular
construction state and by a particular embodiment.

The problem becomes more pronounced at building scale. Navigation through a
floor or room can be evaluated with simplified walls, openings, and obstacles,
whereas placing a connection plate or inserting a pin depends on
millimetre-scale interface geometry. Loading every building object at
interaction-level fidelity is wasteful and can make collision checking and
physics simulation unnecessarily expensive. Retaining only coarse geometry is
equally problematic: it can conceal penetrations, reverse an insertion axis, or
turn an infeasible assembly into an apparently valid one. A useful robot world
must therefore change its working set and geometric fidelity with the task
without losing identity or relationships across scales.

Previous studies have established important parts of this information chain.
BIM-derived robot worlds and task-planning systems have been generated for
painting, prefabricated assembly, and interior wall installation
\cite{kim2021bimrobot,zhu2023ifcsdf,oyediran2025fourdbim}. IFC semantics have
also been linked to action-level robot control through the Robot Construction
Action Node \cite{zhu2024rcan}, while closed-loop digital twins have used
as-built observations and human intervention to adjust BIM-driven construction
workflows \cite{wang2024closeddt}. These studies show that building information
can drive robot planning and simulation. They do not, however, establish a
general criterion for deciding whether a composed construction scene is ready
for a specified robot task before that task is executed.

In parallel, physics-ready asset generation has become considerably more
automated. Xu et al.\ converted construction-equipment CAD data and mechanical
specifications into articulated USD models, using a large language model to
recover missing joint information \cite{xu2024physics}. OpenUSD now provides a
strong technical basis for composing such assets through layers, references,
and payloads \cite{openUSD}. Ren et al.\ have already shown that information
delivery requirements can drive a specification-based IFC-to-USD pipeline
while preserving lifecycle identifiers and model structure
\cite{ren2025ifcusd}. NVIDIA's SimReady framework further organizes
machine-checkable asset requirements into profiles, features, and rules
\cite{nvidiaSimReadyWorkflow}. This progress is essential, but an individually
valid beam, tool, robot, or collider does not imply a valid assembly scene. The
assets may refer to inconsistent coordinate frames, a task may address the
wrong IFC object, the required work-zone payload may be absent, or the
construction state may violate an action precondition. The current SimReady
rules operate at asset and feature level; the official documentation notes
that scene-level specifications and validators are not yet defined
\cite{nvidiaSimReadyFAQ}.

\Cref{fig:scene-gap} summarizes the resulting transition. Asset-oriented
checks can establish that individual models contain valid geometry, collision
shapes, or physical attributes, but the construction task is evaluated over a
composition. Readiness therefore depends on closed cross-asset references,
consistent spatial frames, the active work-zone payload, construction-state
preconditions, and an interaction model whose fidelity matches the intended
action.

\begin{figure*}[t]
\centering
\plannedfigure{38mm}{Figure 1 production placeholder: asset readiness is not
scene readiness}{
Four linked panels: (a) an IFC design model; (b) individually valid robot,
tool, and component assets; (c) composed-scene defects such as a missing
payload, inconsistent frame, wrong task target, and unsupported friction; and
(d) the auditable construction-scene contract required before execution.
Use one IFC GUID and one task reference across the panels to make the failure
chain explicit.}{
Implemented scene fixtures, validator rule identifiers, and deterministic
fault examples. Vector SVG/PDF; no generated simulation imagery.}
\caption{From individually valid digital assets to a task-ready construction
scene. Scene-level failures emerge from composition, construction state, and
task binding even when each asset passes its own checks.}
\label{fig:scene-gap}
\end{figure*}

Construction makes this scene-level gap particularly consequential. Temporary
states such as pre-positioned, aligned, partially inserted, and temporarily
fixed are not normally represented in an as-designed model, although they
determine whether the next robotic action is safe and meaningful. Physical
properties may also be assembled from sources of unequal authority: geometry
and identity may be explicit in IFC, mass may be calculated from volume and
material density, friction may come from a published interval, and an assembly
tolerance may be introduced as a task-specific assumption. If these values are
stored without provenance, a simulator can produce a precise-looking result
whose physical basis cannot be audited. A language-model agent compounds the
problem when it is allowed to invent or modify engineering parameters without
a deterministic acceptance test.

Humanoid robots provide a demanding but useful setting in which to examine
these issues. Human-oriented buildings contain doors, thresholds, narrow work
zones, tools, and operating heights that motivate legged mobility and
whole-body manipulation. This does not make a humanoid universally preferable
to a mobile manipulator; wheeled systems retain clear advantages in payload,
stability, and operation on regular floors. Rather, a humanoid task exposes the
need to connect building-scale circulation, body placement, bimanual reach, and
local contact geometry in one composed scene. The present work therefore uses
light-steel assembly by a Unitree G1-class humanoid as the evaluation context,
while keeping the proposed scene representation independent of a particular
controller. A recent construction-humanoid digital-twin study has already
demonstrated the conversion of BIM geometry and semantic information into a
MuJoCo scene for Unitree G1 task decomposition and simulation
\cite{ye2026humanoiddt}. The question pursued here is consequently not whether
BIM and a humanoid can coexist in a simulator, but whether the resulting
multi-asset scene can be compiled, audited, repaired, and related
quantitatively to task outcome.

This paper develops \systemname{} around the proposition that construction
scene readiness should be compiled and tested explicitly, rather than
discovered only through failed simulation runs. IFC4.3 entities are first
translated into a construction scene intermediate representation that retains
component identity, work zones, assembly interfaces, construction states,
tasks, and property provenance. The representation is then compiled into a
global navigation layer, one or more work-zone payloads, and task-specific
interaction layers. This organization permits OpenUSD to load detailed
geometry only where the active task requires it, while references preserve the
connection to the same engineering objects across layers.

A deterministic contract is evaluated over the composed scene rather than over
isolated assets. Its rules cover coordinate and unit consistency, reference
closure, task-to-interface binding, state preconditions, work-zone activation,
and the availability and provenance of critical physical parameters. When a
rule fails, an agent receives a structured diagnostic and may select only from
a whitelist of repair operations, such as restoring an IFC--USD mapping,
recomputing inertia, or regenerating an allowed payload. The validator, rather
than the agent, decides whether the revised scene is acceptable. Parameters
whose authority or risk exceeds a defined threshold are escalated for human
review.

The evaluation is built around controlled failure rather than a single
successful demonstration. Parameterized IFC4.3 fixtures generate correct
scene descriptions and machine-readable ground truth for connection-plate
positioning, pin insertion, and sequential assembly. Known semantic, spatial,
physical, interface, state, and task defects are then injected to compare
generic asset checks, fixed-rule conversion, unconstrained agentic repair, and
validator-grounded repair. CPU experiments measure detection, repair
reliability, provenance coverage, and the resource effect of task-activated
layers. Subsequent RTX experiments test whether the resulting readiness scores
and critical rule failures explain motion-planning and humanoid assembly
outcomes. In this way, the paper links an information-model claim---that a
scene satisfies an explicit construction contract---to an operational
claim---that the composed world supports the intended robotic task.
